\documentclass[10pt,reqno,a4paper]{amsart}

\usepackage[utf8]{inputenc}
\usepackage[english]{babel}
\usepackage{amsmath,amssymb,amsthm,amsfonts}
\usepackage{mathtools}
\usepackage[margin=25mm]{geometry}
\usepackage{booktabs}
\usepackage{array}
\usepackage{seqsplit}
\usepackage{microtype}
\usepackage[colorlinks=true,linkcolor=blue,urlcolor=blue]{hyperref}
\usepackage[capitalize]{cleveref}

\theoremstyle{plain}
\newtheorem{theorem}{Theorem}[section]
\newtheorem{proposition}[theorem]{Proposition}
\newtheorem{lemma}[theorem]{Lemma}
\theoremstyle{definition}
\newtheorem{definition}[theorem]{Definition}
\theoremstyle{remark}
\newtheorem{remark}[theorem]{Remark}

\newcommand{\R}{\mathbb R}
\newcommand{\Z}{\mathbb Z}
\newcommand{\N}{\mathbb N}
\newcommand{\Pp}{\mathbb P}
\newcommand{\Ee}{\mathbb E}
\newcommand{\Fou}{\mathcal F}
\DeclareMathOperator{\Var}{Var}
\newcommand{\dri}[1]{\lVert #1\rVert_{\mathrm{DRI}}}
% Declaration names are long and TeX cannot break inside \texttt. Turn every
% underscore into a break opportunity, with no hyphen inserted.
\newcommand{\ubreak}{\char`\_\discretionary{}{}{}}
\newcommand{\lean}[1]{{\ttfamily\let\_\ubreak\seqsplit{#1}}}

\title[Two formalized libraries]{Two formalized libraries: a positive-drift
first-passage central limit theorem and the key renewal theorem}
\author{Benny Avelin}
\date{\today}

\raggedbottom

\begin{document}

\begin{abstract}
  The manuscript \emph{Absorption cutoff and stationary singularities for
  rounded Gaussian random dynamical systems} uses two classical results that
  Mathlib does not contain: a central limit theorem for the first passage time
  of a positive-drift random walk, and the key renewal theorem for directly
  Riemann integrable functions on the line. The manuscript cites the literature
  for both. The accompanying Lean development cannot, so it proves them, and
  those proofs appear nowhere in the paper. This note records them at outline
  level and maps each step to the declaration that carries it.
\end{abstract}

\maketitle

\section{Why this note exists}

A formalization cannot cite. Where the manuscript writes ``by
\cite{gut2009stopped}'' or ``by \cite{feller1971vol2}'', the Lean development
has to supply a proof, and in these two cases Mathlib offered no starting
point: it contains no renewal theory whatsoever --- the identifiers
\lean{renewal}, \lean{Blackwell}, \lean{directly Riemann} and \lean{ladder} do
not occur in the library --- and no first-passage asymptotics. Two libraries
were therefore built inside the project.

They are genuine mathematical content of the formalization rather than
bookkeeping, and a reader who wants to know what was actually verified should
not have to reconstruct them from Lean source. That is what this note is for.
It is an outline: complete enough to follow the architecture and check that the
steps compose, and deliberately not a substitute for the source. Two longer
notes in this repository, \texttt{CH3\_FIRST\_PASSAGE\_CLT\_PROOF.tex} and
\texttt{A4G4\_BLACKWELL\_PROOF\_NOTE.tex}, carry the detailed arguments for
\cref{sec:clt} and for the constant identification in \cref{sec:renewal}
respectively.

Throughout, declaration names are given without their \lean{AbsorptionCutoff} or
\lean{AbsorptionCutoff.Renewal} prefix.

\section{The positive-drift first-passage central limit theorem}
\label{sec:clt}

\subsection{Setting}

Let $X_0, X_1, \dots$ be independent, identically distributed, square
integrable real random variables with
\begin{equation*}
  \mu := \Ee X_0 > 0,
  \qquad
  \sigma^2 := \Var X_0 \in (0,\infty),
\end{equation*}
and write
\begin{equation*}
  S_n := \sum_{j<n} X_j,
  \qquad
  M_n := \max_{0\le k\le n} S_k,
  \qquad
  H_u := \inf\{n\ge 0 : S_n \ge u\},
\end{equation*}
for the partial sums (\lean{partialSum}), their running maximum
(\lean{runningMax}) and the first weak ascending passage time above the level
$u$ (\lean{firstPassageTime}, defined as Mathlib's \lean{hittingAfter} of
$[u,\infty)$). The passage time takes values in $\N\cup\{\infty\}$.

\begin{theorem}[Sequence-indexed first-passage profile]\label{thm:clt}
  Let $n_r \to \infty$ be observation times and $u_r$ levels such that
  \begin{equation}\label{eq:clt-centering}
    \frac{u_r - n_r\mu}{\sigma\sqrt{n_r}} \longrightarrow c \in \R .
  \end{equation}
  Then $\Pp(H_{u_r} > n_r) \to \Phi(c)$, where $\Phi$ is the standard Gaussian
  distribution function.
\end{theorem}

\noindent
This is \lean{tendsto\_measureReal\_firstPassageTime\_gt}.

\subsection{Outline}

\emph{Step 1: duality.} Directly from the definitions,
\begin{equation}\label{eq:duality}
  \{H_u > n\} = \{M_n < u\},
\end{equation}
so the passage-time profile is a statement about the running maximum. Everything
below is about $M_n$.

\emph{Step 2: the terminal drawdown by finite reversal.} Put
$D_n := M_n - S_n \ge 0$ (\lean{terminalDrawdown}). Reversing the first $n$
increments --- the permutation \lean{finiteReversal}, i.e.\ \lean{Fin.revPerm}
--- carries $D_n$ to the maximal negative partial sum of the reversed block:
\begin{equation}\label{eq:reversal}
  D_n
  \;\overset{d}{=}\;
  \max_{0\le j\le n}\bigl(-S_j\bigr) .
\end{equation}
The identity is pathwise on a finite block (\lean{backwardPartialSum},
\lean{negativeBackwardRunningMax}, \lean{finiteNegativeRunningMax}), and the two
blocks have the same law because the increments are i.i.d.; only exchangeability
of the finite block is used.

\emph{Step 3: uniform tightness of the drawdown.} By \eqref{eq:reversal} it
suffices to bound $\max_{j\le n}(-S_j)$ uniformly in $n$. Positive drift and the
strong law give $\inf_j S_j > -\infty$ almost surely, so the events
$\{\forall j,\ S_j \ge -m\}$ (\lean{lowerBoundEvent}) increase to a set of full
measure. Hence $\{D_n\}_{n}$ is tight, uniformly in $n$ --- the point being that
the bound does not degrade as $n$ grows.

\emph{Step 4: the central limit theorem for the running maximum.} Write
$M_n = S_n + D_n$. Mathlib's Lindeberg--L\'evy theorem
(\lean{tendstoInDistribution\_inv\_sqrt\_mul\_sum\_sub}) gives
$(S_n - n\mu)/(\sigma\sqrt n) \Rightarrow N(0,1)$, and $D_n/\sqrt n \to 0$ in
probability by Step 3. Slutsky then yields
\begin{equation*}
  \frac{M_n - n\mu}{\sigma\sqrt n} \Longrightarrow N(0,1) .
\end{equation*}
This is the only place a distributional limit is created; everything after it is
deterministic bookkeeping about levels and times.

\emph{Step 5: from the maximum back to the passage time.} By \eqref{eq:duality},
\begin{equation*}
  \Pp(H_{u_r} > n_r)
  =
  \Pp\!\left(
    \frac{M_{n_r} - n_r\mu}{\sigma\sqrt{n_r}}
    <
    \frac{u_r - n_r\mu}{\sigma\sqrt{n_r}}
  \right),
\end{equation*}
and \eqref{eq:clt-centering} together with continuity of $\Phi$ gives
\cref{thm:clt}. Continuity is what makes the strict inequality harmless, and it
is also why the Gaussian limit law being atomless matters here.

\emph{Step 6: canonical and post-floor observation times.} The application needs
a specific time, not an abstract sequence. For a diverging scale $L_r$ and a
Gaussian offset $a$ set
\begin{equation}\label{eq:canonical-time}
  t_r
  :=
  \left\lfloor
    \frac{L_r}{\mu}
    +
    \frac{a\sigma}{\mu^{3/2}}\sqrt{L_r}
    +
    q_r
  \right\rfloor ,
\end{equation}
which is \lean{canonicalTimeArgument} and \lean{canonicalTime}; the variant that
adds the perturbation \emph{after} flooring is \lean{postFloorTime}. Substituting
$n_r = t_r$ and $u_r = L_r$ into \eqref{eq:clt-centering} makes the centering
ratio converge to $-a$, so the limit is $\Phi(-a)$
(\lean{tendsto\_measureReal\_firstPassageTime\_gt\_canonicalTime} and
\lean{tendsto\_measureReal\_firstPassageTime\_gt\_postFloorTime}). The floor
costs at most one unit and is absorbed on the $\sqrt{L_r}$ scale.

\emph{Step 7: an almost surely bounded increasing correction.} The logarithmic
radius of the Gaussian model is not a random walk but a random walk plus a
nonlinear correction. The library therefore proves the profile for
$S_n + E_n$, where $E$ is nonnegative, nondecreasing in $n$, and dominated by an
almost surely finite envelope $E_\infty$; that conjunction is the path event
\lean{correctionGoodEvent}, and the corrected passage time is
\lean{correctedFirstPassageTime}. Monotonicity gives the two-sided sandwich
\begin{equation*}
  H_{u} - \text{(shift)}
  \;\le\;
  H^{E}_{u}
  \;\le\;
  H_{u}
\end{equation*}
at levels displaced by $E_\infty$ (\lean{correctedFirstPassageTime\_sandwich}),
and one squeezes with the auxiliary displacement $B_r := \sqrt{\sqrt{L_r}}$,
chosen because $B_r \to \infty$ --- so the envelope is eventually below it ---
while $B_r/\sqrt{L_r}\to 0$, so the displacement is invisible on the CLT scale.
The result is
\lean{tendsto\_measureReal\_correctedFirstPassageTime\_gt\_postFloorTime}, again
with limit $\Phi(-a)$. This is the form the fixed-width cutoff theorem consumes.

\subsection{Declaration map}

\begin{center}
\renewcommand{\arraystretch}{1.2}
\begin{tabular}{@{}l>{\raggedright\arraybackslash}p{0.60\textwidth}@{}}
  \toprule
  Object or step & Declaration \\
  \midrule
  $S_n$, $M_n$, $D_n$, $H_u$
    & \lean{partialSum}, \lean{runningMax}, \lean{terminalDrawdown},
      \lean{firstPassageTime} \\
  Reversal identity \eqref{eq:reversal}
    & \lean{finiteReversal}, \lean{backwardPartialSum},
      \lean{negativeBackwardRunningMax} \\
  Drawdown tightness
    & \lean{lowerBoundEvent}, \lean{negativeRunningMax} \\
  Abstract profile (\cref{thm:clt})
    & \lean{tendsto\_measureReal\_firstPassageTime\_gt} \\
  Observation times \eqref{eq:canonical-time}
    & \lean{canonicalTimeArgument}, \lean{canonicalTime}, \lean{postFloorTime} \\
  At the canonical time
    & \lean{tendsto\_measureReal\_firstPassageTime\_gt\_canonicalTime} \\
  At the post-floor time
    & \lean{tendsto\_measureReal\_firstPassageTime\_gt\_postFloorTime} \\
  Bounded increasing correction
    & \lean{correctionGoodEvent}, \lean{correctedFirstPassageTime},
      \lean{correctedFirstPassageTime\_sandwich} \\
  Corrected profile
    & \lean{tendsto\_measureReal\_correctedFirstPassageTime\_gt\_postFloorTime} \\
  \bottomrule
\end{tabular}
\end{center}

\noindent
All of the above live in \texttt{AbsorptionCutoff/FirstPassageCLT.lean}, which is a pure
Mathlib leaf: it imports nothing else from the project.

\section{The key renewal theorem}\label{sec:renewal}

\subsection{Setting}

Let $\widehat\mu$ be a nonlattice probability measure on $\R$ with finite second
moment and drift
\begin{equation*}
  \widehat m := \int_\R z\,\widehat\mu(\mathrm dz) \in (0,\infty),
\end{equation*}
and assume a finite exponential moment $\int e^{-\theta z}\,\widehat\mu(\mathrm dz)<1$
for some $\theta>0$. Write $\widehat\mu^{*n}$ for the $n$-fold convolution and
\begin{equation*}
  U := \sum_{n\ge0}\widehat\mu^{*n}
\end{equation*}
for the renewal measure. For $g:\R\to\R$ put
\begin{equation}\label{eq:dri-norm}
  \dri{g} := \sum_{k\in\Z}\ \sup_{y\in[k,k+1]} \lvert g(y)\rvert ,
\end{equation}
the directly-Riemann-integrable norm (\lean{driNorm}).

\begin{theorem}[Key renewal theorem]\label{thm:key-renewal}
  For every continuous $z:\R\to\R$ with $\dri{z}<\infty$,
  \begin{equation*}
    \sum_{n\ge0}\int_\R z(y-s)\,\widehat\mu^{*n}(\mathrm ds)
    \;\longrightarrow\;
    \frac{1}{\widehat m}\int_\R z ,
    \qquad y\to\infty .
  \end{equation*}
\end{theorem}

\noindent
This is \lean{tendsto\_tsum\_integral\_comp\_sub\_of\_driNorm}.

\subsection{Outline}

Mathlib supplies convolution of measures (\lean{Measure.conv}) with
associativity, commutativity, unit laws and the integral formulas, and nothing
else relevant. Everything below is built on that substrate.

\emph{Step 1: the renewal measure is locally finite.} A Chernoff estimate on the
exponential transform gives the explicit left-tail bound
$U(s) \le e^{\theta b}/(1-\int e^{-\theta z}\,\mathrm d\widehat\mu)$ for
$s\subseteq(-\infty,b]$ (\lean{renewalMeasure\_le\_of\_expTransform}), hence
$U(I)<\infty$ for every bounded interval
(\lean{renewalMeasure\_Icc\_lt\_top}). Local finiteness is what makes the
smoothed series meaningful at all.

\emph{Step 2: the d.R.i.\ norm controls the series.} A \emph{uniform} bound
$U([y,y+\ell])\le C$, valid for all $y\in\R$ rather than only $y\to\infty$,
converts \eqref{eq:dri-norm} into the estimate
\begin{equation}\label{eq:dri-control}
  \sum_{n\ge0}\int_\R g(y-s)\,\widehat\mu^{*n}(\mathrm ds)
  \;\le\;
  C\,\dri{g}
  \qquad\text{uniformly in } y
\end{equation}
(\lean{tsum\_lintegral\_comp\_sub\_le}), together with
$\lVert\cdot\rVert_{L^1}\le\dri{\cdot}$
(\lean{lintegral\_le\_driNorm}). Two kernels' series therefore differ by at most
$C\dri{z-w}$ uniformly in $y$
(\lean{norm\_tsum\_integral\_comp\_sub\_sub\_le}). This inequality is the engine
of the whole approximation argument: it means it is enough to prove the limit on
a $\dri{\cdot}$-dense class.

\emph{Step 3: a kernel with compactly supported Fourier transform.} The
Fourier transform of the box $\mathbf 1_{[-1,1]}$ is the sinc function, so the
\emph{square}
\begin{equation*}
  \mathrm{sincSq} := \mathrm{sinc}^2
\end{equation*}
is nonnegative and its Fourier transform is the triangle function, which is
compactly supported (\lean{RenewalKernel.lean}: \lean{sincSq},
\lean{triangle\_eq}, \lean{sincSq\_le}). Compact \emph{spectral} support is
exactly what licenses the Riemann--Lebesgue step without assuming any decay of
the characteristic function of $\widehat\mu$ at infinity --- decay that
nonlatticeness alone does not provide. This is the reason the argument is run
against a bandlimited test kernel rather than directly against $z$.

\emph{Step 4: Abel regularization identifies the constant.} Testing the renewal
series against a bandlimited kernel and inverting the Fourier transform produces
a pole at the origin. The naive symmetric principal value around that pole yields
only \emph{half} of $1/\widehat m$; the missing half is the boundary term that
the one-sided renewal series $\sum_{n\ge0}$ selects. Abel regularization ---
weighting the $n$-th convolution by $r^n$ and letting $r\uparrow1$ --- keeps the
characteristic-function estimates, avoids the improper Dirichlet integral, and
delivers the full constant. The outcome is the reference-kernel limit
\begin{equation*}
  \sum_{n\ge0}\int v(y-z)\,\widehat\mu^{*n}(\mathrm dz)
  \longrightarrow
  \frac{1}{\widehat m}\int v
\end{equation*}
for bounded nonnegative continuous bandlimited $v$
(\lean{tendsto\_tsum\_integral\_comp\_sub}), extended to signed bandlimited
kernels by \lean{tendsto\_of\_tendsto\_sum\_integral\_comp\_sub\_bandlimited}.
\textbf{This is the step where Blackwell's constant is identified}, and it is the
mathematical heart of the library;
\texttt{A4G4\_BLACKWELL\_PROOF\_NOTE.tex} gives it in full, including why the
Laplace--Tauberian alternative does not shorten the work.

\emph{Step 5: discharging the abstract hypotheses.} Reading
$\Fou(\mathrm{sincSq})$ at the origin gives $\int \mathrm{sincSq} = 2$
(\lean{integral\_sincSq}), so Step 4 applies concretely:
$\sum_n\int \mathrm{sincSq}(y-z)\,\widehat\mu^{*n}(\mathrm dz)\to 2/\widehat m$
(\lean{tendsto\_tsum\_integral\_comp\_sub\_sincSq}). The two-sided uniform cell
bound of Step 2 is obtained here as well, by combining the $y\to\infty$ bound
with the Chernoff left-tail bound
(\lean{exists\_bound\_renewalMeasure\_Icc\_of\_expTransform}).

\emph{Step 6: bandlimited approximation in $\dri{\cdot}$.} Let
$K_a := \tfrac a2\,\mathrm{sincSq}(a\,\cdot)$ be the rescaled approximate
identity (\lean{smoothKernel}): it has unit mass, band $[-2a,2a]$ and decay
$4a/(1+a^2t^2)$. For continuous compactly supported $w$ the smoothing
$w \star K_a$ (\lean{smoothed}) is bandlimited at $2a$
(\lean{fourier\_smoothed\_eq\_zero}), converges to $w$ uniformly
(\lean{exists\_forall\_norm\_smoothed\_sub\_le}), and has far field decaying like
$(a\cdot\mathrm{dist}^2)^{-1}$ (\lean{norm\_smoothed\_le\_of\_support}) --- fast
enough for the far cells of \eqref{eq:dri-norm} to be summable. The two regimes
combine into $\dri{w - w\star K_a}\to0$
(\lean{exists\_driNorm\_smoothed\_sub\_lt}).

\emph{Step 7: conclusion.} A continuous kernel of finite d.R.i.\ norm is
$\dri{\cdot}$-approximable by continuous compactly supported ones
(\lean{exists\_hasCompactSupport\_driNorm\_sub\_lt}), those by bandlimited ones
(Step 6), and the limit is known for bandlimited kernels (Steps 4--5). Since by
\eqref{eq:dri-control} the series depend on the kernel $\dri{\cdot}$-Lipschitzly
and uniformly in $y$, the limit passes to $z$, giving \cref{thm:key-renewal}.

\begin{remark}[Why continuity, and not almost-everywhere continuity]
  The classical statement assumes only that $z$ is almost everywhere continuous.
  This route needs genuine continuity, because the approximation is carried out
  in $\dri{\cdot}$ and indicators are \emph{not} $\dri{\cdot}$-approximable by
  continuous functions --- the supremum over a cell straddling a jump does not
  shrink. The gap is real at the level of the abstract theorem and costs nothing
  in the application: the manuscript states \lean{lem:nd-gaussian-renewal} with
  the continuity hypothesis, and \lean{prop:nd-forcing-admissibility} proves that
  the only forcing the lemma is ever applied to is continuous in the radial
  level.
\end{remark}

\subsection{Declaration map}

\begin{center}
\renewcommand{\arraystretch}{1.2}
\begin{tabular}{@{}l>{\raggedright\arraybackslash}p{0.60\textwidth}@{}}
  \toprule
  Object or step & Declaration \\
  \midrule
  Renewal measure, Chernoff tail
    & \lean{renewalMeasure\_le\_of\_expTransform},
      \lean{renewalMeasure\_Icc\_lt\_top} \\
  d.R.i.\ norm \eqref{eq:dri-norm} and \eqref{eq:dri-control}
    & \lean{driNorm}, \lean{tsum\_lintegral\_comp\_sub\_le},
      \lean{lintegral\_le\_driNorm} \\
  Uniform kernel-difference bound
    & \lean{norm\_tsum\_integral\_comp\_sub\_sub\_le} \\
  Uniform cell bound
    & \lean{exists\_bound\_renewalMeasure\_Icc\_of\_expTransform} \\
  The $\mathrm{sinc}^2$ kernel
    & \lean{sincSq}, \lean{triangle\_eq}, \lean{integral\_sincSq} \\
  Constant identification (Abel)
    & \lean{tendsto\_tsum\_integral\_comp\_sub} \\
  Signed bandlimited kernels
    & \lean{tendsto\_of\_tendsto\_sum\_integral\_comp\_sub\_bandlimited} \\
  Concrete limit at $\mathrm{sinc}^2$
    & \lean{tendsto\_tsum\_integral\_comp\_sub\_sincSq} \\
  Approximate identity and smoothing
    & \lean{smoothKernel}, \lean{smoothed}, \lean{fourier\_smoothed\_eq\_zero} \\
  Bandlimited $\dri{\cdot}$-approximation
    & \lean{exists\_driNorm\_smoothed\_sub\_lt},
      \lean{exists\_hasCompactSupport\_driNorm\_sub\_lt} \\
  Key renewal theorem (\cref{thm:key-renewal})
    & \lean{tendsto\_tsum\_integral\_comp\_sub\_of\_driNorm} \\
  \bottomrule
\end{tabular}
\end{center}

\noindent
These live in \texttt{AbsorptionCutoff/Supercritical/}: \texttt{Renewal.lean} (renewal
measure, d.R.i.\ norm, cell estimates), \texttt{RenewalKernel.lean} (the
$\mathrm{sinc}^2$ kernel), \texttt{RenewalAbel.lean} (the Abel argument),
\texttt{RenewalSinc.lean} (discharging the abstract hypotheses) and
\texttt{RenewalApprox.lean} (bandlimited approximation and the theorem). The
split across five modules is for build time, not mathematics.

\subsection{What the manuscript uses}

\cref{thm:key-renewal} is instantiated at the Cram\'er-tilted increment law
$\widehat\mu_{A,N}$ to give the manuscript's
\lean{lem:nd-gaussian-renewal}, whose four analytic hypotheses --- nonlattice,
square integrable, positive drift, finite exponential moment --- are each
established for that law in \texttt{StationaryEquation.lean}. The angular part
of the manuscript's renewal equation requires no measure theory on the sphere:
the product form of the tilted kernel makes the angular integral factor out as a
constant, leaving the same scalar convolution for every test function.

\begin{thebibliography}{9}
\bibitem{feller1971vol2}
  W.~Feller,
  \emph{An Introduction to Probability Theory and Its Applications}, vol.~2,
  2nd ed., Wiley, 1971. \S XI.1 and \S XI.9.
\bibitem{gut2009stopped}
  A.~Gut,
  \emph{Stopped Random Walks: Limit Theorems and Applications},
  2nd ed., Springer, 2009. Theorem 3.5.1; Theorem 6.6 and Remark 6.4.
\end{thebibliography}

\end{document}
